\documentclass[a4paper]{amsbook}

\usepackage{color}
\usepackage[colorlinks=true, citecolor=red]{hyperref}  %needs to come before amsrefs package to work with \cite{}*{}
\usepackage{amssymb,latexsym}
\usepackage[alphabetic]{amsrefs}
\usepackage[mathscr]{eucal}
\usepackage{pdfsync}
\usepackage{graphicx}
\usepackage{epstopdf}
\DeclareGraphicsRule{.tif}{png}{.png}{`convert #1 `dirname #1`/`basename #1 .tif`.png}
\usepackage[all]{xy}
\usepackage{titletoc}
%\usepackage[show]{ed}
%\usepackage{txfonts}   %  \fint  for integral average

%\usepackage{showkeys}  %options notref  notcite 

%\input{rgb.tex}

%\textcolor{DarkRed}{text in DarkRed}


\CompileMatrices

\theoremstyle{plain}
\newtheorem{theorem}{Theorem}[section]
\newtheorem{lemma}[theorem]{Lemma}
\newtheorem{proposition}[theorem]{Proposition}
\newtheorem{corollary}[theorem]{Corollary}

\theoremstyle{definition}
\newtheorem{definition}[theorem]{Definition}%[section]


\theoremstyle{remark}
\newtheorem{defn}[theorem]{Definition} 
\newtheorem{example}[theorem]{Example}
\newtheorem{remark}[theorem]{Remark}
\newtheorem{discussion}[theorem]{Discussion}

%\numberwithin{equation}{section}

%\newcommand{\abs}[1]{\lvert#1\rvert}

\DeclareMathOperator{\Lip}{Lip}
\DeclareMathOperator{\dist}{dist}
\DeclareMathOperator{\expected}{\mathbb{E}}
\DeclareMathOperator{\prb}{Prob}
\DeclareMathOperator{\spt}{spt}
\DeclareMathOperator{\lap}{\Delta}
\DeclareMathOperator{\Tr}{Tr}

\raggedbottom

\addtolength{\textwidth}{2cm}
\addtolength{\oddsidemargin}{-1cm}
\addtolength{\evensidemargin}{-1cm} 
\addtolength{\topmargin}{-1cm} 
\addtolength{\textheight}{2cm}

\newcommand{\side}[1]{\footnote{\emph{Draft Comment}:  #1}}
%\newcommand{\side}[1]{} 

\def\Xint#1{\mathchoice
   {\XXint\displaystyle\textstyle{#1}}%
   {\XXint\textstyle\scriptstyle{#1}}%
   {\XXint\scriptstyle\scriptscriptstyle{#1}}%
   {\XXint\scriptscriptstyle\scriptscriptstyle{#1}}%
   \!\int}
\def\XXint#1#2#3{{\setbox0=\hbox{$#1{#2#3}{\int}$}
     \vcenter{\hbox{$#2#3$}}\kern-.5\wd0}}
\def\ddashint{\Xint=}
\def\dashint{\Xint-}




%%% BEGIN DOCUMENT
\begin{document}

\title{Index Theory}  
\author{John E.\ Hutchinson}
\date{\today}

 \begin{abstract}
  I follow Gilkey's book
 \end{abstract}

\maketitle



\tableofcontents

\chapter{Pseudo Differential Operators}
%%%%%%%%%%%%%%%%%%%%%%%%%%%%%%%%
%%%%%%%%%%%%%%%%%%%%%%%%%%%%%%%%
%%%%%%%%%%%%%%%%%%%%%%%%%%%%%%%%
%%%%%%%%%%%%%%%%%%%%%%%%%%%%%%%%
\section{\textbf{Fourier Transforms, Schwartz class, Sobolev Spaces}} 

%%%%%%%%%%%%%%%%%%%%%%%%%%
 \subsection{Notation}  We work in $  \mathbb{R}^m$.  For $ \alpha_i \geq 0 $ 
  \begin{gather*}
  \alpha  = (\alpha_1,\dots, \alpha_m )   , \quad 
  | \alpha |   = \alpha_1+ \dots + \alpha _m, \\ 
  \alpha !  = \alpha _1 ! \cdot \dots\cdot  \alpha _m ! , \quad 
  x^ \alpha  = x_1^{\alpha _1} \cdot \dots \cdot x_m ^ {\alpha _m} \\
  \partial^ \alpha = \partial_x^ \alpha 
     = \left( \frac{\partial}{\partial x_1}\right)^{\alpha_1}\cdots   \left( \frac{\partial}{\partial x_m}\right)^{\alpha_m} , \quad 
D^ \alpha  = D^ \alpha _x = (-i)^{| \alpha | }  d^ \alpha _x .
\end{gather*}
 
 \emph{Taylor's formula}: For $ f \in C^\infty $,
 \[
 f(x) = \sum_{|\alpha   |\leq n} \frac{1}{ \alpha !} \partial^ \alpha f(x_0) (x-x_0)^ \alpha + O(|x-x_0|).
 \]
 
 \emph{Schwartz class}: 
 \[
  \mathcal{S} = \{ f: \mathbb{R}^m \to \mathbb{C} \mid f \in C^\infty , \  \forall n  \forall \beta \ | D^\beta  f(x) | \leq C_{n, \beta }(1+ |x|)^{-n}\},
  \]
  where $ C_{n, \beta } $ depend on $ f $.
  
  To simplify the FT formulae, Lebesgue measure on $ \mathbb{R}^m $ is normalised by $ (2\pi)^{-m/2} $, so that
  \[
  \int e^{-\frac{1}{2} |x|^2} dx   = 1.
  \]
  
 \emph{Convolution}: 
  Define
  \[
  (f*g)(x) := \int f(x-y)\, g(y)\, dy = \int f(y)\, g(x-y)\,  dy.
  \]
  This defines an \emph{associative and commutative multiplication}.
  
  \begin{proposition}  If $ f \in \mathcal{S} $, $ \int f = 1 $, $ f_\epsilon (x) := \epsilon ^{-n}f \left(\frac{x}{\epsilon } \right) $, 
 and $g  \in \mathcal{S} $, then  
\[
f_\epsilon * g \to g \quad \text{uniformly  as } \epsilon \to 0 .
\]  
\end{proposition}

One can use this approach to show any \emph{$ g\in L^p $ can be approximated in the $ L^p $ norm by a $ C^\infty $ function of compact support}.  
  
  
%%%%%%%%%%%%%%%%%%%%%%%%%%%%%%%%%%%%%%%  
\subsection{Definition of FT on $ \mathcal{S}$ and basic properties.}
 Define 
 \[
 \hat{f}( \xi ) = \int e^{-ix\cdot \xi } f(x) \, dx.
 \]
 
 It follows, using integration by parts and dominated convergence for the first line, that %
 \footnote{Typo in book}
 \footnote{Where ``$ \leftrightarrow $'' means the right side is the FT of the left side.}
\begin{gather}
\notag  D^\alpha_x f (x)  \leftrightarrow \xi^ \alpha \hat{f}( \xi ), 
                                    \quad  x^ \alpha f(x) \leftrightarrow D^ \alpha_\xi \hat{f} ( \xi ),\\
 f(x-x_0)  \leftrightarrow e^{-ix_0\cdot \xi} \hat{f}(\xi), 
                   \quad e^{ ix\cdot \xi_0 } f(x) \leftrightarrow \hat{f}(\xi-\xi_0). \label{tm}
\end{gather} 
 In particular, it follows $ \hat{f} \in \mathcal{S} $.

Calculation shows
\[
e^{-\frac{1}{2} |x|^2} \leftrightarrow e^{-\frac{1}{2} |\xi|^2} .
\]

\begin{proposition}  $ f \leftrightarrow \hat{f} $ is a bijection on $ \mathcal{S} $.  In fact
\[
f(x) = \int e^{ix\cdot \xi } \hat{f} (\xi)\, d\xi = \hat{\hat{f}} (-x) =: Tf(x).
\]
\end{proposition}

\begin{proof}
Suppose $ f \in \mathcal{S}$.  

First suppose $ f(0) = 0 $.  Use
\[
f(x) = \int_0^1 f(tx)\, dt 
\quad 
\text{to get} \quad 
f(x) = \sum_j x_j h_j(x)  \text{ where } h_j \in \mathcal{S} ,
\]
and so%
\footnote{missing ``$ - $'' in book.} 
\[
\hat{f}  = -i\sum_j \frac{\partial \hat{h}_j}{\partial \xi_j}  , \quad 
\therefore Tf(0) = \int \hat{f}  = 0 = f(0).
\]

For general $ f $ subtract a multiple of $ e^{-\frac{1}{2} |x|^ 2 }$ and note that $ T $ is the identity on $ e^{-\frac{1}{2} |x|^ 2 }$.  It follows $ Tf(0) = f(0) $ for any $ f \in \mathcal{S} $.

To get equality at any $ x \in \mathbb{R}^m $ use \eqref{tm}.
\end{proof} 

We also have
\begin{gather*}
f * g \leftrightarrow \hat{f}\cdot \hat{g}, \quad f \cdot g \leftrightarrow \hat{f}* \hat{g} \quad \forall f,g\in \mathcal{S}, \\
(f,g) = (\hat{f}, \hat{g} ) \quad \forall f,g \in L^2(\mathbb{R}^m ).
\end{gather*}
The last gives a unitary map on $ L^2( \mathbb{R}^m)$ and is known as \emph{Plancherel's theorem}.

  
 
%%%%%%%%%%%%%%%%%%%%%%%%%%%%%%%%%%%%%%%%%%
\subsection{Sobolev Spaces} \mbox{}

\emph{Definition}:  For $ f \in \mathcal{S} $ and $ s \in \mathbb{R} $  let
\[
|f|^2_s = \int(1+|\xi|^2 )^s |\hat{f}(\xi)|\, d\xi .
\]
Note $ (1+|\xi|^2)^s $ and $ (1+|\xi|)^{2s} $ give equivalent norms.
$ \mathcal{H}_s(\mathbb{R}^m ) $ is the completion of $ \mathcal{S} $ in this norm.

\medskip
\emph{Equivalent spaces}: $ \mathcal{H}_0\approx L^2 $ by Plancherel's theorem.
$ \mathcal{H}_s\approx L^2 $ if the measure $ (1+|\xi|^2)^{s/2} $ is used.


\medskip
\emph{Connection with derivatives}:  $ s $ counts the number of derivatives in the sense 
\[
\sum_{| \alpha |\leq n} \int |D^ \alpha f(x) |^2 = \sum_{| \alpha |\leq n}  \int \left| \xi ^ \alpha \hat{f}( \xi) \right|^2 
         \sim \| f \|^2_s,
         \]
  since
  \[
  c_1(1+|\xi|^2)^n  \leq \sum_{| \alpha |\leq n} |\xi^ \alpha |^2 \leq c_2(1+|\xi|^2)^n .
  \]
  
  
  \medskip  \emph{Continuous map}: The map $ D^ \alpha \mathcal{H} _s \to \mathcal{H} _{s- | \alpha | }$ is continuous.
  
\medskip \emph{Higher derivatives imply continuity}:  If $ s > k + m/2 $ and $  \mathcal{H}_s  \subset C^k $ and
\[
|f|_{\infty,k} \leq C |f|_s .
\]

\medskip \emph{Dense subset}:  $ C_0^\infty  $ is dense in $ \mathcal{H}_s $.

\medskip \emph{Rellich's lemma}:  Suppose $ (f_n)_{n\geq 1} \in \mathcal{S} $ with uniform compact support, and uniform $ \mathcal{H}_s $ bound.  Then if $ t < s $, some subsequence converges in $ \mathcal{H}_t $.


\medskip \emph{Natural dual pairing}: The map $ (f,g) \mapsto \int f\overline{g} $ extends to a map $ \mathcal{H}_s \times \mathcal{H}_{-s} \to \mathbb{C} $ which identifies $ \mathcal{H}_{-s} $ with $ \mathcal{H}_s^* $.  In fact
\begin{align*}
|(f,g)| & \leq |f|_s |g|_{-s},\\
|f|_s &= \sup \{ |(f,g)| : |g|_{-s} = 1 \}.
\end{align*}
For the proof of (1) note 
\[
|(f,g)| = |(\hat{f},\hat{g}) = \int \hat{f}(\xi) (1+|\xi|^2)^{s/2}\,  \overline{\hat{g}}(\xi) (1+|\xi|^2)^{-s/2},
\]
and use Cauchy Schwartz.  For (2) choose $ g $ so $ \hat{g} = \hat{f}(1+|\xi |^1)^s $ and normalise.

\medskip \emph{Peter Paul principle}: If $ s>t>u $ and $ \epsilon > 0 $ then 
\[
(1+|\xi|)^{2t} \leq \epsilon (1+|\xi|)^{2s} +  C(\epsilon) (1+|\xi|)^{2s},
\]
and so  
\[
|f|_t \leq \epsilon |f|_s + C(\epsilon) |f|_u.
\]

\medskip \emph{Vector-valued case}:  All the preceding generalises readily to functions with values in a $ k $ dimensional space.







%%%%%%%%%%%%%%%%%%%%%%%%%%%%%%%%
%%%%%%%%%%%%%%%%%%%%%%%%%%%%%%%%
%%%%%%%%%%%%%%%%%%%%%%%%%%%%%%%%
%%%%%%%%%%%%%%%%%%%%%%%%%%%%%%%%


%\begin{thebibliography}{9} 

%

%\end{thebibliography}

\end{document} 

